\begin{myexercise}
Verwenden Sie das Cheatsheet, um eine Tabelle mit folgenden Kolonnen zu erstellen:
\begin{table}[H]
\centering
\begin{tabular}{c|c|c|c}
\texttt{Vorname}&\texttt{Nachname}&\texttt{Klasse}&\texttt{Geburtsdatum}\\\midrule\midrule
...&...&...&...\\
\end{tabular}
\caption{Tabellenname: \texttt{SchuelerInnen}}
\end{table}
Fügen Sie danach 2 Einträge von Schülerinnen Ihrer Klasse in die Tabelle ein.
\end{myexercise}




\subsection*{Lernziele}
\begin{todolist}
\item Ich kann Tabellen erstellen, löschen und verändern, indem ich (unter anderem) die Befehle \lstinline|CREATE TABLE|, \lstinline|DROP TABLE| und \lstinline|ALTER TABLE| verwenden.
\item Ich kann Einträge (Zeilen) in einer existierenden Tabelle hinzufügen, verändern oder löschen, indem ich die Befehle \lstinline|INSERT INTO|, \lstinline|UPDATE| und \lstinline|DELETE FROM| verwende.
\item Ich kann einfache SQL-Abfragen formulieren mit Befehlen wie \lstinline|SELECT|, \lstinline|WHERE|,  \lstinline|ORDER BY|, \lstinline|LIMIT| und weiteren Befehlen (s. Cheatsheet für vollständige Liste aller zu lernenden SQL-Befehle)
\end{todolist}.